\subsection{A sparse Li criterion for every number field}

The preceding prime-sum test loses off-line zeros when their real parts
approach an unattained supremum. The same block argument gives a full
equivalence for a different arithmetic sequence. For a number field $K$
set $w_\rho=1-1/\rho$ and define its Li coefficients by
\[
 \lambda_n(K)=\lim_{T\to\infty}\sum_{|\operatorname{Im}\rho|\le T}
            (1-w_\rho^n),\qquad n\ge1.
\]
The sums are real by conjugation. Equivalently, if $\xi_K$ is the entire
completed Dedekind zeta function with its poles removed, then
$\lambda_n(K)=\frac{1}{(n-1)!}(d/ds)^n[s^{n-1}\log\xi_K(s)]_{s=1}$.

\paragraph{Corollary.}
For every number field $K$, critical-line RH for $\zeta_K$ is equivalent
to $\lambda_n(K)\ge0$ for all sufficiently large $n\in\mathcal E$.
The same equivalence holds if the fraction of negative coefficients
in the $j$th block tends to zero. The index set samples at most
$O(1+\log\log N)$ coefficients of index at most $N$.

\paragraph{Proof.}
On RH, $|w_\rho|=1$ and every real part $1-\operatorname{Re}(w_\rho^n)$
is nonnegative. If RH fails, symmetry about the critical line gives
a zero with real part below $1/2$, hence $|w_\rho|>1$ since
$|w_\rho|^2=1+(1-2\operatorname{Re}\rho)/|\rho|^2$. The numbers
$w_\rho$ tend to $1$ at increasing height. Consequently the largest
modulus $R>1$ is attained at finitely many zeros, say $m$ with
multiplicity, and all other moduli are at most $r$ for some $1<r<R$.
For all these other zeros Taylor's formula on the segment from $1$ to
$w_\rho$ gives, for $n\ge2$,
\[
 \left|1-w_\rho^n-\frac n\rho\right|
 \le \frac{n(n-1)}{2|\rho|^2}r^{n-2}.
\]
The sum of $|\rho|^{-2}$ is finite by the zero count. The symmetric sum
of $1/\rho$ exists since paired conjugates have sum
$2\operatorname{Re}\rho/|\rho|^2$. It follows that
\[
 R^{-n}\lambda_n(K)
 =-\sum_{|w_\rho|=R}(w_\rho/R)^n+o(1).
\]
Apply the same-box argument to the arguments of the $m$ outermost
points. The real parts of their normalized powers can all be made
arbitrarily close to $1$ in every sufficiently large block, and the
real part of their sum never exceeds $m$. Thus
\[
 \lim_{j\to\infty}\min_{1\le k\le j}
 R^{-kM_j}\lambda_{kM_j}(K)=-m.
\]
Taking the phase tolerance fixed and small shows that a positive
fraction of every large block has negative coefficients if RH fails.
Finally, counting blocks with $M_j\le N$ gives the claimed sparsity.

Li's number-field criterion and the dominant exponential terms in
its failure case are classical. The additional conclusion is the
prescribed sparse set of indices, its robustness to a vanishing
fraction of exceptions in each block, and the normalized block limit.
In fact the proof permits any sequence of positive integer dilations
$M_j\to\infty$ with block lengths tending to infinity. For
$K=\mathbb Q(\sqrt[3]{2})$, for example, proving the sampled
inequalities would establish RH for that field, which includes a
nonabelian factor. We have derived the equivalence. An arithmetic proof
of the inequalities for this field remains open.
See E.~Bombieri and J.~C.~Lagarias,
\href{https://websites.umich.edu/lagarias/doc/bombieri.pdf}{\emph{Complements to Li's
criterion for the Riemann hypothesis}}, Theorem~1, and
J.~C.~Lagarias,
\href{https://doi.org/10.5802/aif.2311}{\emph{Li coefficients for automorphic
$L$-functions}}, especially the introduction on the Dedekind case.
